\pdfoutput=1
%% bare_jrnl_compsoc.tex
%% V1.4b
%% 2015/08/26
%% by Michael Shell
%% See:
%% http://www.michaelshell.org/
%% for current contact information.
%%
%% This is a skeleton file demonstrating the use of IEEEtran.cls
%% (requires IEEEtran.cls version 1.8b or later) with an IEEE
%% Computer Society journal paper.
%%
%% Support sites:
%% http://www.michaelshell.org/tex/ieeetran/
%% http://www.ctan.org/pkg/ieeetran
%% and
%% http://www.ieee.org/

%%*************************************************************************
%% Legal Notice:
%% This code is offered as-is without any warranty either expressed or
%% implied; without even the implied warranty of MERCHANTABILITY or
%% FITNESS FOR A PARTICULAR PURPOSE! 
%% User assumes all risk.
%% In no event shall the IEEE or any contributor to this code be liable for
%% any damages or losses, including, but not limited to, incidental,
%% consequential, or any other damages, resulting from the use or misuse
%% of any information contained here.
%%
%% All comments are the opinions of their respective authors and are not
%% necessarily endorsed by the IEEE.
%%
%% This work is distributed under the LaTeX Project Public License (LPPL)
%% ( http://www.latex-project.org/ ) version 1.3, and may be freely used,
%% distributed and modified. A copy of the LPPL, version 1.3, is included
%% in the base LaTeX documentation of all distributions of LaTeX released
%% 2003/12/01 or later.
%% Retain all contribution notices and credits.
%% ** Modified files should be clearly indicated as such, including  **
%% ** renaming them and changing author support contact information. **
%%*************************************************************************


% *** Authors should verify (and, if needed, correct) their LaTeX system  ***
% *** with the testflow diagnostic prior to trusting their LaTeX platform ***
% *** with production work. The IEEE's font choices and paper sizes can   ***
% *** trigger bugs that do not appear when using other class files.       ***
% The testflow support page is at:
% http://www.michaelshell.org/tex/testflow/


\documentclass[10pt,journal,compsoc]{IEEEtran}
%
% If IEEEtran.cls has not been installed into the LaTeX system files,
% manually specify the path to it like:
% \documentclass[10pt,journal,compsoc]{../sty/IEEEtran}





% Some very useful LaTeX packages include:
% (uncomment the ones you want to load)


% *** MISC UTILITY PACKAGES ***
%
%\usepackage{ifpdf}
% Heiko Oberdiek's ifpdf.sty is very useful if you need conditional
% compilation based on whether the output is pdf or dvi.
% usage:
% \ifpdf
%   % pdf code
% \else
%   % dvi code
% \fi
% The latest version of ifpdf.sty can be obtained from:
% http://www.ctan.org/pkg/ifpdf
% Also, note that IEEEtran.cls V1.7 and later provides a builtin
% \ifCLASSINFOpdf conditional that works the same way.
% When switching from latex to pdflatex and vice-versa, the compiler may
% have to be run twice to clear warning/error messages.






% *** CITATION PACKAGES ***
%
\ifCLASSOPTIONcompsoc
  % IEEE Computer Society needs nocompress option
  % requires cite.sty v4.0 or later (November 2003)
  \usepackage[nocompress]{cite}
\else
  % normal IEEE
  \usepackage{cite}
\fi
% cite.sty was written by Donald Arseneau
% V1.6 and later of IEEEtran pre-defines the format of the cite.sty package
% \cite{} output to follow that of the IEEE. Loading the cite package will
% result in citation numbers being automatically sorted and properly
% "compressed/ranged". e.g., [1], [9], [2], [7], [5], [6] without using
% cite.sty will become [1], [2], [5]--[7], [9] using cite.sty. cite.sty's
% \cite will automatically add leading space, if needed. Use cite.sty's
% noadjust option (cite.sty V3.8 and later) if you want to turn this off
% such as if a citation ever needs to be enclosed in parenthesis.
% cite.sty is already installed on most LaTeX systems. Be sure and use
% version 5.0 (2009-03-20) and later if using hyperref.sty.
% The latest version can be obtained at:
% http://www.ctan.org/pkg/cite
% The documentation is contained in the cite.sty file itself.
%
% Note that some packages require special options to format as the Computer
% Society requires. In particular, Computer Society  papers do not use
% compressed citation ranges as is done in typical IEEE papers
% (e.g., [1]-[4]). Instead, they list every citation separately in order
% (e.g., [1], [2], [3], [4]). To get the latter we need to load the cite
% package with the nocompress option which is supported by cite.sty v4.0
% and later. Note also the use of a CLASSOPTION conditional provided by
% IEEEtran.cls V1.7 and later.





% *** GRAPHICS RELATED PACKAGES ***
%

\ifCLASSINFOpdf
  \usepackage[pdftex]{graphicx}
  \usepackage[export]{adjustbox}
  \usepackage{epstopdf}
  % declare the path(s) where your graphic files are
  \graphicspath{{../output/figures/}}
  % and their extensions so you won't have to specify these with
  % every instance of \includegraphics
  \DeclareGraphicsExtensions{.pdf,.png,.eps}
\else
  % or other class option (dvipsone, dvipdf, if not using dvips). graphicx
  % will default to the driver specified in the system graphics.cfg if no
  % driver is specified.
  \usepackage[dvips]{graphicx}
  % declare the path(s) where your graphic files are
  \graphicspath{{../output/figures/}}
  % and their extensions so you won't have to specify these with
  % every instance of \includegraphics
  \DeclareGraphicsExtensions{.eps}
\fi
% graphicx was written by David Carlisle and Sebastian Rahtz. It is
% required if you want graphics, photos, etc. graphicx.sty is already
% installed on most LaTeX systems. The latest version and documentation
% can be obtained at: 
% http://www.ctan.org/pkg/graphicx
% Another good source of documentation is "Using Imported Graphics in
% LaTeX2e" by Keith Reckdahl which can be found at:
% http://www.ctan.org/pkg/epslatex
%
% latex, and pdflatex in dvi mode, support graphics in encapsulated
% postscript (.eps) format. pdflatex in pdf mode supports graphics
% in .pdf, .jpeg, .png and .mps (metapost) formats. Users should ensure
% that all non-photo figures use a vector format (.eps, .pdf, .mps) and
% not a bitmapped formats (.jpeg, .png). The IEEE frowns on bitmapped formats
% which can result in "jaggedy"/blurry rendering of lines and letters as
% well as large increases in file sizes.
%
% You can find documentation about the pdfTeX application at:
% http://www.tug.org/applications/pdftex






% *** MATH PACKAGES ***
%
%\usepackage{amsmath}
% A popular package from the American Mathematical Society that provides
% many useful and powerful commands for dealing with mathematics.
%
% Note that the amsmath package sets \interdisplaylinepenalty to 10000
% thus preventing page breaks from occurring within multiline equations. Use:
%\interdisplaylinepenalty=2500
% after loading amsmath to restore such page breaks as IEEEtran.cls normally
% does. amsmath.sty is already installed on most LaTeX systems. The latest
% version and documentation can be obtained at:
% http://www.ctan.org/pkg/amsmath





% *** SPECIALIZED LIST PACKAGES ***
%
%\usepackage{algorithmic}
% algorithmic.sty was written by Peter Williams and Rogerio Brito.
% This package provides an algorithmic environment fo describing algorithms.
% You can use the algorithmic environment in-text or within a figure
% environment to provide for a floating algorithm. Do NOT use the algorithm
% floating environment provided by algorithm.sty (by the same authors) or
% algorithm2e.sty (by Christophe Fiorio) as the IEEE does not use dedicated
% algorithm float types and packages that provide these will not provide
% correct IEEE style captions. The latest version and documentation of
% algorithmic.sty can be obtained at:
% http://www.ctan.org/pkg/algorithms
% Also of interest may be the (relatively newer and more customizable)
% algorithmicx.sty package by Szasz Janos:
% http://www.ctan.org/pkg/algorithmicx




% *** ALIGNMENT PACKAGES ***
%
%\usepackage{array}
% Frank Mittelbach's and David Carlisle's array.sty patches and improves
% the standard LaTeX2e array and tabular environments to provide better
% appearance and additional user controls. As the default LaTeX2e table
% generation code is lacking to the point of almost being broken with
% respect to the quality of the end results, all users are strongly
% advised to use an enhanced (at the very least that provided by array.sty)
% set of table tools. array.sty is already installed on most systems. The
% latest version and documentation can be obtained at:
% http://www.ctan.org/pkg/array


% IEEEtran contains the IEEEeqnarray family of commands that can be used to
% generate multiline equations as well as matrices, tables, etc., of high
% quality.




% *** SUBFIGURE PACKAGES ***
%\ifCLASSOPTIONcompsoc
%  \usepackage[caption=false,font=footnotesize,labelfont=sf,textfont=sf]{subfig}
%\else
%  \usepackage[caption=false,font=footnotesize]{subfig}
%\fi
% subfig.sty, written by Steven Douglas Cochran, is the modern replacement
% for subfigure.sty, the latter of which is no longer maintained and is
% incompatible with some LaTeX packages including fixltx2e. However,
% subfig.sty requires and automatically loads Axel Sommerfeldt's caption.sty
% which will override IEEEtran.cls' handling of captions and this will result
% in non-IEEE style figure/table captions. To prevent this problem, be sure
% and invoke subfig.sty's "caption=false" package option (available since
% subfig.sty version 1.3, 2005/06/28) as this is will preserve IEEEtran.cls
% handling of captions.
% Note that the Computer Society format requires a sans serif font rather
% than the serif font used in traditional IEEE formatting and thus the need
% to invoke different subfig.sty package options depending on whether
% compsoc mode has been enabled.
%
% The latest version and documentation of subfig.sty can be obtained at:
% http://www.ctan.org/pkg/subfig




% *** FLOAT PACKAGES ***
%
%\usepackage{fixltx2e}
% fixltx2e, the successor to the earlier fix2col.sty, was written by
% Frank Mittelbach and David Carlisle. This package corrects a few problems
% in the LaTeX2e kernel, the most notable of which is that in current
% LaTeX2e releases, the ordering of single and double column floats is not
% guaranteed to be preserved. Thus, an unpatched LaTeX2e can allow a
% single column figure to be placed prior to an earlier double column
% figure.
% Be aware that LaTeX2e kernels dated 2015 and later have fixltx2e.sty's
% corrections already built into the system in which case a warning will
% be issued if an attempt is made to load fixltx2e.sty as it is no longer
% needed.
% The latest version and documentation can be found at:
% http://www.ctan.org/pkg/fixltx2e


%\usepackage{stfloats}
% stfloats.sty was written by Sigitas Tolusis. This package gives LaTeX2e
% the ability to do double column floats at the bottom of the page as well
% as the top. (e.g., "\begin{figure*}[!b]" is not normally possible in
% LaTeX2e). It also provides a command:
%\fnbelowfloat
% to enable the placement of footnotes below bottom floats (the standard
% LaTeX2e kernel puts them above bottom floats). This is an invasive package
% which rewrites many portions of the LaTeX2e float routines. It may not work
% with other packages that modify the LaTeX2e float routines. The latest
% version and documentation can be obtained at:
% http://www.ctan.org/pkg/stfloats
% Do not use the stfloats baselinefloat ability as the IEEE does not allow
% \baselineskip to stretch. Authors submitting work to the IEEE should note
% that the IEEE rarely uses double column equations and that authors should try
% to avoid such use. Do not be tempted to use the cuted.sty or midfloat.sty
% packages (also by Sigitas Tolusis) as the IEEE does not format its papers in
% such ways.
% Do not attempt to use stfloats with fixltx2e as they are incompatible.
% Instead, use Morten Hogholm'a dblfloatfix which combines the features
% of both fixltx2e and stfloats:
%
% \usepackage{dblfloatfix}
% The latest version can be found at:
% http://www.ctan.org/pkg/dblfloatfix




%\ifCLASSOPTIONcaptionsoff
%  \usepackage[nomarkers]{endfloat}
% \let\MYoriglatexcaption\caption
% \renewcommand{\caption}[2][\relax]{\MYoriglatexcaption[#2]{#2}}
%\fi
% endfloat.sty was written by James Darrell McCauley, Jeff Goldberg and 
% Axel Sommerfeldt. This package may be useful when used in conjunction with 
% IEEEtran.cls'  captionsoff option. Some IEEE journals/societies require that
% submissions have lists of figures/tables at the end of the paper and that
% figures/tables without any captions are placed on a page by themselves at
% the end of the document. If needed, the draftcls IEEEtran class option or
% \CLASSINPUTbaselinestretch interface can be used to increase the line
% spacing as well. Be sure and use the nomarkers option of endfloat to
% prevent endfloat from "marking" where the figures would have been placed
% in the text. The two hack lines of code above are a slight modification of
% that suggested by in the endfloat docs (section 8.4.1) to ensure that
% the full captions always appear in the list of figures/tables - even if
% the user used the short optional argument of \caption[]{}.
% IEEE papers do not typically make use of \caption[]'s optional argument,
% so this should not be an issue. A similar trick can be used to disable
% captions of packages such as subfig.sty that lack options to turn off
% the subcaptions:
% For subfig.sty:
% \let\MYorigsubfloat\subfloat
% \renewcommand{\subfloat}[2][\relax]{\MYorigsubfloat[]{#2}}
% However, the above trick will not work if both optional arguments of
% the \subfloat command are used. Furthermore, there needs to be a
% description of each subfigure *somewhere* and endfloat does not add
% subfigure captions to its list of figures. Thus, the best approach is to
% avoid the use of subfigure captions (many IEEE journals avoid them anyway)
% and instead reference/explain all the subfigures within the main caption.
% The latest version of endfloat.sty and its documentation can obtained at:
% http://www.ctan.org/pkg/endfloat
%
% The IEEEtran \ifCLASSOPTIONcaptionsoff conditional can also be used
% later in the document, say, to conditionally put the References on a 
% page by themselves.




% *** PDF, URL AND HYPERLINK PACKAGES ***
%
\usepackage{url}
% url.sty was written by Donald Arseneau. It provides better support for
% handling and breaking URLs. url.sty is already installed on most LaTeX
% systems. The latest version and documentation can be obtained at:
% http://www.ctan.org/pkg/url
% Basically, \url{my_url_here}.





% *** Do not adjust lengths that control margins, column widths, etc. ***
% *** Do not use packages that alter fonts (such as pslatex).         ***
% There should be no need to do such things with IEEEtran.cls V1.6 and later.
% (Unless specifically asked to do so by the journal or conference you plan
% to submit to, of course. )


% correct bad hyphenation here
\hyphenation{op-tical net-works semi-conduc-tor}

% specific packages
\usepackage{hyperref}
\usepackage[capitalize]{cleveref}
\usepackage{xcolor}
\usepackage{gensymb}
\usepackage{ifthen}
\usepackage[percent]{overpic}
\usepackage{stackengine}
\usepackage{longtable}
\usepackage{multirow}
\usepackage{pdflscape}
\usepackage{geometry}

\geometry{
	a4paper,
	left=2mm,
	right=2mm,
	top=10mm,
	bottom=10mm
}

\newcommand{\cf}{\emph{cf}~}
\newcommand{\eg}{\emph{e.g.}~}
\newcommand{\ie}{\emph{i.e.}~}
\newcommand{\vs}{\emph{vs}~}
\newcommand{\AuthorName}{Tomasz Gałaj}

\newcommand{\supplemental}[1]{}%{#1}
%\newcommand{\supplemental}[1]{#1}

\newlength{\mytablewidth}
\setlength{\mytablewidth}{\columnwidth}
\newenvironment{mytabular}{%
  \setlength{\tabcolsep}{0cm}
  \renewcommand\arraystretch{0}
  \tabular%
}{%
  \endtabular
}

\newcommand{\myfigure}[2]{%
\begin{figure*}\begin{center}#1\end{center}#2\end{figure*}%
}

\newlength\lunderset
\newlength\rulethick
\lunderset=2.0pt\relax
\rulethick=0.8pt\relax
\def\stackalignment{l}
\newcommand\nunderline[2][1]{\setbox0=\hbox{#2}%
  \stackunder[#1\lunderset-\rulethick]{\strut#2}{\rule{\wd0}{\rulethick}}}

\title{Supplementary material}
\date{2022\\ July}
\author{Tomasz Ga{\l}aj\\ Adam Wojciechowski}

\makeatletter
\def\@maketitle{%
	%\newpage
	%\null
	%\vskip 2em%
	\begin{center}%
		\let \footnote \thanks
		{\LARGE \@title \par}%
		\vskip 1.5em%
		{\large
			\lineskip .5em%
			\begin{tabular}[t]{c}%
				\@author
			\end{tabular}\par}%
		%\vskip 1em%
		%{\large \@date}%
	\end{center}%
	%\par
	%\vskip 1.5em
}
\makeatother

\begin{document}
\onecolumn
\maketitle

\begin{table*}[ht!]
\renewcommand{\arraystretch}{1.3}
\caption{{\bf Qualitative comparison}. Summary of the qualitive evaluation of
the 8 clear sky models. The precomputation time and memory complexity for the
Haber and Nishita96 models is for a single Sun zenith angle $\theta_s$ only --
the Bruneton and Elek models precomputations are for $n$ such angles and the
Nishita93 precomputations are independent of $\theta_s$.}
\label{fig:qualitative_comparison}
\centering
\begin{tabular}{|r||c|c|c|c|c|c|c|c|c|}
\hline
\bfseries & \bfseries Supported & \bfseries Aerial & \bfseries Sunset &
\bfseries Scattering & \bfseries Precompute & \bfseries Precompute &
\bfseries Render & \bfseries RMSE & \bfseries Delta E \\
\bfseries Model & \bfseries viewpoints & \bfseries perspective & \bfseries
sunrise & \bfseries orders & \bfseries time & \bfseries memory & \bfseries time
& \bfseries $mW/(m^2.sr.nm)$ & \bfseries error \\
\hline
Nishita93 & all & yes & yes & 1 & $O(n^3)$ & $O(n^2)$ & $O(n)$ &
\input{../output/figures/error_total_nishita93.txt} &
\input{../output/lab_diffs/lab_diff_nishita93_avg_delta_e.txt} \\
\hline
Nishita96 & in atmosphere & yes & yes & 2 & $O(n^3)$ & $O(n^3)$ & $O(n)$ &
\input{../output/figures/error_total_nishita96.txt} &
\input{../output/lab_diffs/lab_diff_nishita96_avg_delta_e.txt} \\
\hline
%Preetham & ground only & yes & no & 2 & $0$ & $0$ & $O(1)$ &
%\input{../output/figures/error_total_preetham.txt} \\
%\hline
O'Neal & all & yes & yes & $1$ & $0$ & $0$ & $O(n)$ &
\input{../output/figures/error_total_oneal.txt} &
\input{../output/lab_diffs/lab_diff_oneal_avg_delta_e.txt} \\
\hline
Haber & ground only & yes & yes & all & $O(n^6)$ & $O(n^3)$ & $O(n^2)$ &
\input{../output/figures/error_total_haber.txt} &
\input{../output/lab_diffs/lab_diff_haber_avg_delta_e.txt} \\
\hline
Bruneton & all & yes & yes & all & $O(n^6)$ & $O(n^4)$ & $O(1)$ &
\input{../output/figures/error_total_bruneton.txt} &
\input{../output/lab_diffs/lab_diff_bruneton_avg_delta_e.txt} \\
\hline
Elek & all & yes & yes & all & $O(n^6)$ & $O(n^4)$ & $O(1)$ &
\input{../output/figures/error_total_elek.txt} &
\input{../output/lab_diffs/lab_diff_elek_avg_delta_e.txt} \\
\hline
Hosek & ground only & no & no & all & $0$ & $0$ & $O(1)$ &
\input{../output/figures/error_total_hosek.txt} \\
\hline
Trapezoidal & all & yes & yes & 1 & $0$ & $0$ & $O(n^2)$ &
\input{../output/figures/error_total_trapezoidal.txt} &
\input{../output/lab_diffs/lab_diff_trapezoidal_avg_delta_e.txt} \\
\hline
Spline & all & yes & yes & 1 & n/a & n/a & $O(n)$ &
\input{../output/figures/error_total_spline.txt} &
\input{../output/lab_diffs/lab_diff_spline_avg_delta_e.txt} \\
\hline
Taylor & all & yes & yes & 1 & n/a & n/a & $O(n)$ &
\input{../output/figures/error_total_taylor.txt} &
\input{../output/lab_diffs/lab_diff_taylor_avg_delta_e.txt} \\
\hline
\end{tabular}
\end{table*}


%\newpage
\def \my_img_width {0.16} % 1.98/12 = 0.165
\newcommand{\modelrow}[2]{
\rotatebox{90}{\footnotesize #2 /
\input{../output/figures/sza_#2.txt}\unskip\degree} &
\includegraphics[width=\my_img_width \mytablewidth]{#1_#2_nishita93} &
\includegraphics[width=\my_img_width \mytablewidth]{#1_#2_nishita96} &
\includegraphics[width=\my_img_width \mytablewidth]{#1_#2_preetham} &
\includegraphics[width=\my_img_width \mytablewidth]{#1_#2_oneal} &
\includegraphics[width=\my_img_width \mytablewidth]{#1_#2_haber} &
\includegraphics[width=\my_img_width \mytablewidth]{#1_#2_bruneton} &
\includegraphics[width=\my_img_width \mytablewidth]{#1_#2_elek} &
\includegraphics[width=\my_img_width \mytablewidth]{#1_#2_hosek} &
\includegraphics[width=\my_img_width \mytablewidth]{#1_#2_trapezoidal} &
\includegraphics[width=\my_img_width \mytablewidth]{#1_#2_spline} &
\includegraphics[width=\my_img_width \mytablewidth]{#1_#2_taylor} &
\includegraphics[width=\my_img_width \mytablewidth]{#1_#2_libradtran} &
\ifthenelse{\equal{#2}{06h00} \OR \equal{#2}{07h00} \OR \equal{#2}{08h00}}
{\raisebox{0.091\mytablewidth}{\makebox[\my_img_width \mytablewidth]{\footnotesize n/a}}}
{\includegraphics[width=\my_img_width \mytablewidth]{#1_#2_measurements}}}

\newcommand{\modeltable}[3]{
\setlength{\tabcolsep}{0cm}
\renewcommand\arraystretch{0}
\LTcapwidth=\textwidth
\begin{longtable}{p{0.04\mytablewidth}*{14}{p{\my_img_width \mytablewidth}}}
\caption*{#2}\\
&
\footnotesize\bfseries Nishita93 &
\footnotesize\bfseries Nishita96 &
\footnotesize\bfseries Preetham &
\footnotesize\bfseries O'Neal &
\footnotesize\bfseries Haber &
\footnotesize\bfseries Bruneton &
\footnotesize\bfseries Elek &
\footnotesize\bfseries Hosek &
\footnotesize\bfseries Trapezoidal &
\footnotesize\bfseries\leavevmode\color{violet} Spline &
\footnotesize\bfseries\leavevmode\color{orange} Taylor &
\footnotesize\bfseries\leavevmode\color{blue} libRadtran &
\footnotesize\bfseries\leavevmode\color{red} Measurements\vspace{1mm}\\
\modelrow{#1}{06h00}\\
\supplemental{\modelrow{#1}{07h00}\\}
\supplemental{\modelrow{#1}{08h00}\\}
\supplemental{\modelrow{#1}{09h30}\\}
\supplemental{\modelrow{#1}{09h45}\\}
\supplemental{\modelrow{#1}{10h00}\\}
\modelrow{#1}{10h15}\\
\supplemental{\modelrow{#1}{10h30}\\}
\supplemental{\modelrow{#1}{10h45}\\}
\supplemental{\modelrow{#1}{11h00}\\}
\modelrow{#1}{11h15}\\
\supplemental{\modelrow{#1}{11h30}\\}
\supplemental{\modelrow{#1}{11h45}\\}
\supplemental{\modelrow{#1}{12h00}\\}
\supplemental{\modelrow{#1}{12h15}\\}
\supplemental{\modelrow{#1}{12h30}\\}
\supplemental{\modelrow{#1}{12h45}\\}
\supplemental{\modelrow{#1}{13h00}\\}
\modelrow{#1}{13h15}\\
\supplemental{\modelrow{#1}{13h30}\\}
\multicolumn{15}{c}{#3}
\end{longtable}
}

%\myfigure{\modeltable{image_full_spectral}}
\modeltable{image_full_spectral}{%
{\bf Rendering}. Fisheye skydome rendering of the spectral radiance
obtained with each model, convolved with the CIE color matching functions,
converted from XYZ to linear sRGB, and tone mapped with $1-e^{-kL}$, for several
time of day / sun zenith angle values (the red cross indicates the sun
direction). The measurements are interpolated using bicubic spherical
interpolation before rendering. Compare with Fig. 13
in\cite{Kider2014}.\label{fig:rendering}}{\ignorespaces}


\newpage
\modeltable{absolute_luminance}{%
{\bf Absolute luminance}. The absolute luminance in $cd.m^{-2}$
obtained with each model, using the same color scale as in Fig.~7
in~\cite{Zotti2007}. The measurements are interpolated using bicubic spherical
interpolation before being converted to luminance
values.\label{fig:absolute_luminance}}{
{\includegraphics[width=1.5\mytablewidth]{scale_caption}\vspace{-0.5cm}}}

\newpage
\modeltable{relative_luminance}{%
{\bf Relative luminance}. The luminance obtained with each model, in
percentage of the zenith luminance, and using the same color scale as in Fig.~6
in~\cite{Zotti2007}. The measurements are interpolated using bicubic spherical
interpolation before being converted to luminance
values.\label{fig:relative_luminance}}{\includegraphics[width=1.5\mytablewidth]{scale_caption}\vspace{-0.5cm}}


\newpage
\modeltable{chromaticity}{%
{\bf Chromaticity}. The rg-chromaticity, $(r,g,b)/\mathrm{max}(r,g,b)$,
obtained with each model. The measurements are interpolated using bicubic
spherical interpolation before being converted to chromaticity
values.\label{fig:chromaticity}}{\ignorespaces}


\newpage
\newcommand{\relativeerrorcell}[2]{%
\begin{overpic}[width=\my_img_width \mytablewidth]{../output/lab_diffs/lab_diff_#2_#1}
\put(5,5){\footnotesize\textcolor{white}{\input{../output/lab_diffs/lab_diff_#2_#1_delta_e.txt}}}
\end{overpic}}

\newcommand{\relativeerrorrow}[1]{
\rotatebox{90}{\footnotesize #1 / 
\input{../output/figures/sza_#1.txt}\unskip\degree} &
\relativeerrorcell{nishita93}{#1} &
\relativeerrorcell{nishita96}{#1} &
%\relativeerrorcell{preetham}{#1} &
\relativeerrorcell{oneal}{#1} &
\relativeerrorcell{haber}{#1} &
\relativeerrorcell{bruneton}{#1} &
\relativeerrorcell{elek}{#1} &
%\relativeerrorcell{hosek}{#1} &
\relativeerrorcell{trapezoidal}{#1} &
\relativeerrorcell{spline}{#1} &
\relativeerrorcell{taylor}{#1} &
\relativeerrorcell{libradtran}{#1}}

\newcommand{\relativeerrortable}[2]{
\begin{longtable}{p{0.04\mytablewidth}*{13}{p{\my_img_width \mytablewidth}}}
\caption*{#1}\\
&
\footnotesize\bfseries Nishita93 &
\footnotesize\bfseries Nishita96 &
%\footnotesize\bfseries Preetham &
\footnotesize\bfseries O'Neal &
\footnotesize\bfseries Haber &
\footnotesize\bfseries Bruneton &
\footnotesize\bfseries Elek &
%\footnotesize\bfseries Hosek &
\footnotesize\bfseries Trapezoidal &
\footnotesize\bfseries Spline &
\footnotesize\bfseries Taylor &
\footnotesize\bfseries\leavevmode\color{blue} libRadtran\vspace{1mm}\\
\supplemental{\relativeerrorrow{09h30}\\}
\supplemental{\relativeerrorrow{09h45}\\}
\supplemental{\relativeerrorrow{10h00}\\}
\relativeerrorrow{10h15}\\
\supplemental{\relativeerrorrow{10h30}\\}
\supplemental{\relativeerrorrow{10h45}\\}
\supplemental{\relativeerrorrow{11h00}\\}
\relativeerrorrow{11h15}\\
\supplemental{\relativeerrorrow{11h30}\\}
\supplemental{\relativeerrorrow{11h45}\\}
\supplemental{\relativeerrorrow{12h00}\\}
\supplemental{\relativeerrorrow{12h15}\\}
\supplemental{\relativeerrorrow{12h30}\\}
\supplemental{\relativeerrorrow{12h45}\\}
\supplemental{\relativeerrorrow{13h00}\\}
\relativeerrorrow{13h15}\\
\supplemental{\relativeerrorrow{13h30}\\}
\vspace{0.2cm}\\
\multicolumn{14}{c}{#2}
\end{longtable}}

\relativeerrortable{%
{\bf Delta E error}. The Delta E difference (per pixel) computed in CIE L*a*b* space.\label{fig:relative_lab_error}}
{\includegraphics[width=1.5\mytablewidth, height=0.3cm, frame]{../output/lab_diffs/hot_colormap}}


\newpage
\renewcommand{\relativeerrorcell}[2]{%
\begin{overpic}[width=\my_img_width \mytablewidth]{relative_error_#2_#1}
\put(5,5){\footnotesize\input{../output/figures/error_#2_#1.txt}}
\end{overpic}}

\renewcommand{\relativeerrorrow}[1]{
\rotatebox{90}{\footnotesize #1 / 
\input{../output/figures/sza_#1.txt}\unskip\degree} &
\relativeerrorcell{nishita93}{#1} &
\relativeerrorcell{nishita96}{#1} &
\relativeerrorcell{preetham}{#1} &
\relativeerrorcell{oneal}{#1} &
\relativeerrorcell{haber}{#1} &
\relativeerrorcell{bruneton}{#1} &
\relativeerrorcell{elek}{#1} &
\relativeerrorcell{hosek}{#1} &
\relativeerrorcell{trapezoidal}{#1} &
\relativeerrorcell{spline}{#1} &
\relativeerrorcell{taylor}{#1} &
\relativeerrorcell{libradtran}{#1}}

\renewcommand{\relativeerrortable}[2]{
\begin{longtable}{p{0.04\mytablewidth}*{13}{p{\my_img_width \mytablewidth}}}
\caption*{#1}\\
&
\footnotesize\bfseries Nishita93 &
\footnotesize\bfseries Nishita96 &
\footnotesize\bfseries Preetham &
\footnotesize\bfseries O'Neal &
\footnotesize\bfseries Haber &
\footnotesize\bfseries Bruneton &
\footnotesize\bfseries Elek &
\footnotesize\bfseries Hosek &
\footnotesize\bfseries Trapezoidal &
\footnotesize\bfseries Spline &
\footnotesize\bfseries Taylor &
\footnotesize\bfseries\leavevmode\color{blue} libRadtran\vspace{1mm}\\
\supplemental{\relativeerrorrow{09h30}\\}
\supplemental{\relativeerrorrow{09h45}\\}
\supplemental{\relativeerrorrow{10h00}\\}
\relativeerrorrow{10h15}\\
\supplemental{\relativeerrorrow{10h30}\\}
\supplemental{\relativeerrorrow{10h45}\\}
\supplemental{\relativeerrorrow{11h00}\\}
\relativeerrorrow{11h15}\\
\supplemental{\relativeerrorrow{11h30}\\}
\supplemental{\relativeerrorrow{11h45}\\}
\supplemental{\relativeerrorrow{12h00}\\}
\supplemental{\relativeerrorrow{12h15}\\}
\supplemental{\relativeerrorrow{12h30}\\}
\supplemental{\relativeerrorrow{12h45}\\}
\supplemental{\relativeerrorrow{13h00}\\}
\relativeerrorrow{13h15}\\
\supplemental{\relativeerrorrow{13h30}\\}
\multicolumn{14}{c}{#2}
\end{longtable}}

\relativeerrortable{%
{\bf Relative error}. The relative error compared with the
measurements, in \% and using the same color scale as in~\cite{Kider2014},
computed at the 81 sampling points (and summed equally over the common range
supported by all models, i.e. between $360$ and $720\,nm$), and then
interpolated with spherical bicubic interpolation. The bottom left number is the
RMSE in $mW/(m^2.sr.nm)$ (computed over the 81 sampling points and all the
wavelengths between $360$ and $720\,nm$).\label{fig:relative_error}}
{\includegraphics[width=1.5\mytablewidth]{error_caption}\vspace{-0.5cm}}


\newpage
\def \myuserstudyimgwidth {0.45}

\newcommand{\userstudyrow}[1]{
\includegraphics[width=\myuserstudyimgwidth\mytablewidth]{sky_image_#1_nishita93}&
\includegraphics[width=\myuserstudyimgwidth\mytablewidth]{sky_image_#1_nishita96}&
%\includegraphics[width=\myuserstudyimgwidth\mytablewidth]{sky_image_#1_preetham}&
\includegraphics[width=\myuserstudyimgwidth\mytablewidth]{sky_image_#1_oneal}&
\includegraphics[width=\myuserstudyimgwidth\mytablewidth]{sky_image_#1_haber}&
\includegraphics[width=\myuserstudyimgwidth\mytablewidth]{sky_image_#1_bruneton}&
\includegraphics[width=\myuserstudyimgwidth\mytablewidth]{sky_image_#1_elek}
%\includegraphics[width=\myuserstudyimgwidth\mytablewidth]{sky_image_#1_hosek}
}

\newcommand{\userstudy}{
\setlength{\tabcolsep}{0cm}
\renewcommand\arraystretch{0}
\LTcapwidth=\textwidth
\begin{longtable}{*{6}{p{\myuserstudyimgwidth\mytablewidth}}}
\caption*{{\bf Examplary rendered viewpoints from the ground level}. The images were generated for each model for different daytimes (i.e. late afternoon - top, morning - bottom).\label{fig:userstudy}}\\

\footnotesize\bfseries Nishita93 &
\footnotesize\bfseries Nishita96 &
%\footnotesize\bfseries Preetham &
\footnotesize\bfseries O'Neal &
\footnotesize\bfseries Haber &
\footnotesize\bfseries Bruneton &
\footnotesize\bfseries Elek 
%\footnotesize\bfseries Hosek 
\vspace{1mm} \\ 
\userstudyrow{sunrise}\\
\userstudyrow{morning}
\end{longtable}
}

%%%%%%%%%%%%%%%%%%%

\newcommand{\userstudyroww}[1]{
	\includegraphics[width=\myuserstudyimgwidth\mytablewidth]{sky_image_#1_trapezoidal}&
	\includegraphics[width=\myuserstudyimgwidth\mytablewidth]{sky_image_#1_spline} &
	\includegraphics[width=\myuserstudyimgwidth\mytablewidth]{sky_image_#1_taylor} 
}

\newcommand{\userstudyy}{
	\setlength{\tabcolsep}{0cm}
	\renewcommand\arraystretch{0}
	\LTcapwidth=\textwidth
	\begin{longtable}{*{6}{p{\myuserstudyimgwidth\mytablewidth}}}
		\caption*{}\\
		\footnotesize\bfseries Trapezoidal &
		\footnotesize\bfseries Spline &
		\footnotesize\bfseries Taylor
		\vspace{1mm} \\ 
		\userstudyroww{sunrise}\\
		\userstudyroww{morning}
	\end{longtable}
}

\begin{landscape}
	\parbox[c][\textwidth][s]{\linewidth}{%
		%\vfill
		
		\userstudy
		\userstudyy
		
		\vfill}
\end{landscape}



\bibliographystyle{unsrt}
\bibliography{references}

\end{document}